\begin{abstract}

The differing representation spaces required for visual understanding and generation pose a challenge in unifying them within the autoregressive paradigm of large language models. A vision tokenizer trained for reconstruction excels at capturing low-level visual appearance, making it well-suited for visual generation but lacking high-level semantic representations for understanding tasks. Conversely, a vision encoder trained via contrastive learning aligns well with language but struggles to decode back into the pixel space for generation tasks.
To bridge this gap, we propose \textbf{\ours}, a method that unifies representations for both understanding and generation within a single tokenizer.
However, directly integrating reconstruction and semantic objectives creates conflicts, leading to degraded performance in both reconstruction fidelity and semantic accuracy.
Instead of forcing a single codebook to capture both visual appearance and semantics, \ours{} disentangles them by introducing separate codebooks for high-level semantics and low-level visual details.
As a result, \ours{} achieves 0.25 rFID and 82.0\% zero-shot accuracy on ImageNet, and demonstrates strong effectiveness in downstream MLLM tasks for both understanding and generation.
Specifically, our method surpasses VILA-U by 5.8 points on average across ten visual understanding benchmarks and delivers a 13\% improvement on GenAI-Bench.
Notably, incorporating dual visual tokens outperforms using a single token type on both understanding and generation tasks.
We hope our research offers a new perspective on leveraging dual visual vocabularies for building unified vision–language models.
Project page is available \href{https://songweii.github.io/dualtoken-project-page/}{here}.

\end{abstract}
